\documentclass[11pt, ukrainian]{article}
\usepackage[a4paper]{geometry}
\setlength\parindent{0mm}
\geometry{top=1.1cm,bottom=1.1cm,left=1.0cm,right=1.0cm}
\pagestyle{empty}
\usepackage[T2A]{fontenc}
\usepackage[utf8]{inputenc}
\usepackage[ukrainian]{babel}
\usepackage{graphicx}
\usepackage{caption}
\DeclareCaptionFormat{nocaption}{}
\captionsetup{format=nocaption,aboveskip=0pt,belowskip=0pt}
\usepackage[Export]{adjustbox} 
\adjustboxset{max size={0.9\linewidth}{0.9\paperheight}}
\begin{document}
%begin
%\begin {Huge}
{{27.11.2025 Контрольна робота \No 2 (частина 2), варіант  \No \textbf { 1 }\par}}

%\begin{center}Частина 2
%\end{center}
%\vspace{10mm}
%\end {Huge}

{
\begin {large}
Якість коду та економність обчислень оцінюється по всім завданням.

%\vspace{5mm}

%beginitem
1. Написати функцію, що повертає найбільше число з списку цілих, у сімковому запису якого нема цифри 2. За відсутності має повертати None. Використовувати додаткові структури даних (у тому числі рядки) заборонено.

%enditem

%beginitem
2. 
Написати функцію, що повертає \textbf{новий список} з записаною в нього послідовністю чисел
$$\scalebox{1.2}{$C_{n+k}^{k+1}$}, \quad k=0,1,\ldots,n-2$$
Використовувати **, вкладені цикли та виклики функцій (крім методів класу list) \textbf{заборонено}.

%enditem
\end{large}
%end
\newpage
%begin
%\begin {Huge}
{{27.11.2025 Контрольна робота \No 2 (частина 2), варіант  \No \textbf { 2 }\par}}

%\begin{center}Частина 2
%\end{center}
%\vspace{10mm}
%\end {Huge}

{
\begin {large}
Якість коду та економність обчислень оцінюється по всім завданням.

%\vspace{5mm}

%beginitem
1. Написати функцію, що для цілого числа знаходить простий дільник, який входить у розклад числа в найбільшому степені. Якщо таких кілька, то повернути найбільший. \\ Наприклад, для числа $2^{9}{\cdot}3^5{\cdot}7^9{\cdot}11$ функція має повернути $7$.

%enditem

%beginitem
2. 
Написати функцію, що повертає \textbf{новий список} з записаною в нього послідовністю чисел
$$\scalebox{1.2}{$C_{2n+k}^{n+k}$}, \quad k=0,1,\ldots,n$$
Використовувати **, вкладені цикли та виклики функцій (крім методів класу list) \textbf{заборонено}.

%enditem
\end{large}
%end
\newpage
%begin
%\begin {Huge}
{{27.11.2025 Контрольна робота \No 2 (частина 2), варіант  \No \textbf { 3 }\par}}

%\begin{center}Частина 2
%\end{center}
%\vspace{10mm}
%\end {Huge}

{
\begin {large}
Якість коду та економність обчислень оцінюється по всім завданням.

%\vspace{5mm}

%beginitem
1. Написати функцію, що повертає найбільше число з списку цілих, що має не більше трьох різних простих дільників. За відсутності має повертати None. Використовувати додаткові структури даних (у тому числі рядки) заборонено.

%enditem

%beginitem
2. 
Написати функцію, що повертає \textbf{новий список} з записаною в нього послідовністю чисел
$$\scalebox{1.2}{$(-1)^{k}C_{2n+k+1}^{n+k}$}, \quad k=1,\ldots,n$$
Використовувати **, вкладені цикли та виклики функцій (крім методів класу list) \textbf{заборонено}.

%enditem
\end{large}
%end
\newpage
%begin
%\begin {Huge}
{{27.11.2025 Контрольна робота \No 2 (частина 2), варіант  \No \textbf { 4 }\par}}

%\begin{center}Частина 2
%\end{center}
%\vspace{10mm}
%\end {Huge}

{
\begin {large}
Якість коду та економність обчислень оцінюється по всім завданням.

%\vspace{5mm}

%beginitem
1. Написати функцію, що повертає число з списку цілих, що має найбільшу кількість різних простих дільників. Якщо таких чисел кілька, то повернути найменше. Використовувати додаткові структури даних (у тому числі рядки) заборонено.

%enditem

%beginitem
2. 
Написати функцію, що повертає \textbf{новий список} з записаною в нього послідовністю чисел
$$\scalebox{1.2}{$C_{2k+1}^{k}$}, \quad k=2,\ldots,n$$
Використовувати **, вкладені цикли та виклики функцій (крім методів класу list) \textbf{заборонено}.

%enditem
\end{large}
%end
\newpage
%begin
%\begin {Huge}
{{27.11.2025 Контрольна робота \No 2 (частина 2), варіант  \No \textbf { 5 }\par}}

%\begin{center}Частина 2
%\end{center}
%\vspace{10mm}
%\end {Huge}

{
\begin {large}
Якість коду та економність обчислень оцінюється по всім завданням.

%\vspace{5mm}

%beginitem
1. Написати функцію, що для цілого числа знаходить кількість простих дільників, які входять у розклад числа в найбільшому степені. \\ Наприклад, для числа $3^5{\cdot}7^9{\cdot}11{\cdot}23^{9}$ таких дільників два – $23$ та $7$, тому функція має повернути $2$.

%enditem

%beginitem
2. 
Написати функцію, що повертає \textbf{новий список} з записаною в нього послідовністю чисел
$$\scalebox{1.2}{$(-1)^{k}C_{2n+k}^{n+k}$}, \quad k=0,1,\ldots,n$$
Використовувати **, вкладені цикли та виклики функцій (крім методів класу list) \textbf{заборонено}.

%enditem
\end{large}
%end
\newpage
%begin
%\begin {Huge}
{{27.11.2025 Контрольна робота \No 2 (частина 2), варіант  \No \textbf { 6 }\par}}

%\begin{center}Частина 2
%\end{center}
%\vspace{10mm}
%\end {Huge}

{
\begin {large}
Якість коду та економність обчислень оцінюється по всім завданням.

%\vspace{5mm}

%beginitem
1. Написати функцію, що для списку цілих знаходить довжину найдовшого відрізку, в якому всі сусідні числа відрізняються якнайменш на 2. Використовувати додаткові структури даних (у тому числі рядки) заборонено.

%enditem

%beginitem
2. 
Написати функцію, що повертає \textbf{новий список} з записаною в нього послідовністю чисел
$$\scalebox{1.2}{$C_{2n+k}^{n+k}$}, \quad k=1,\ldots,n$$
Використовувати **, вкладені цикли та виклики функцій (крім методів класу list) \textbf{заборонено}.

%enditem
\end{large}
%end
\newpage
%begin
%\begin {Huge}
{{27.11.2025 Контрольна робота \No 2 (частина 2), варіант  \No \textbf { 7 }\par}}

%\begin{center}Частина 2
%\end{center}
%\vspace{10mm}
%\end {Huge}

{
\begin {large}
Якість коду та економність обчислень оцінюється по всім завданням.

%\vspace{5mm}

%beginitem
1. Написати функцію, що для інтервалу $[a, b)$ знаходить найближчу пару чисел з цього інтервалу, що кожне з них має непарну кількість різних простих дільників. Якщо таких пар кілька, то цікавить пара з найменшим добутком.

%enditem

%beginitem
2. 
Написати функцію, що повертає \textbf{новий список} з записаною в нього послідовністю чисел
$$\scalebox{1.2}{$C_{n+k}^{2k}$}, \quad k=1,\ldots,n$$
Використовувати **, вкладені цикли та виклики функцій (крім методів класу list) \textbf{заборонено}.

%enditem
\end{large}
%end
\newpage
%begin
%\begin {Huge}
{{27.11.2025 Контрольна робота \No 2 (частина 2), варіант  \No \textbf { 8 }\par}}

%\begin{center}Частина 2
%\end{center}
%\vspace{10mm}
%\end {Huge}

{
\begin {large}
Якість коду та економність обчислень оцінюється по всім завданням.

%\vspace{5mm}

%beginitem
1. Написати функцію, що для списку цілих знаходить найдовший відрізок, в якому всі сусідні числа різні. Якщо таких відрізків кілька, то повернути останній (як індекс початку та довжину). Використовувати додаткові структури даних (у тому числі рядки) заборонено.

%enditem

%beginitem
2. 
Написати функцію, що повертає \textbf{новий список} з записаною в нього послідовністю чисел
$$\scalebox{1.2}{$(-1)^{k}C_{n+2k}^{k}$}, \quad k=0,1,\ldots,n-2$$
Використовувати **, вкладені цикли та виклики функцій (крім методів класу list) \textbf{заборонено}.

%enditem
\end{large}
%end
\newpage
%begin
%\begin {Huge}
{{27.11.2025 Контрольна робота \No 2 (частина 2), варіант  \No \textbf { 9 }\par}}

%\begin{center}Частина 2
%\end{center}
%\vspace{10mm}
%\end {Huge}

{
\begin {large}
Якість коду та економність обчислень оцінюється по всім завданням.

%\vspace{5mm}

%beginitem
1. Написати функцію, що для інтервалу $[a, b)$ знаходить найближчу пару чисел з цього інтервалу, що кожне з них має парну кількість різних простих дільників. Якщо таких пар кілька, то цікавить пара з найменшим добутком.

%enditem

%beginitem
2. 
Написати функцію, що повертає \textbf{новий список} з записаною в нього послідовністю чисел
$$\scalebox{1.2}{$C_{3k-1}^{k}$}, \quad k=2,\ldots,n$$
Використовувати **, вкладені цикли та виклики функцій (крім методів класу list) \textbf{заборонено}.

%enditem
\end{large}
%end
\newpage
%begin
%\begin {Huge}
{{27.11.2025 Контрольна робота \No 2 (частина 2), варіант  \No \textbf { 10 }\par}}

%\begin{center}Частина 2
%\end{center}
%\vspace{10mm}
%\end {Huge}

{
\begin {large}
Якість коду та економність обчислень оцінюється по всім завданням.

%\vspace{5mm}

%beginitem
1. Написати функцію, що для інтервалу $[a, b)$ знаходить знаходить найбільший степінь числа 3, на який ділиться хоча б одне ціле число з цього інтервалу.

%enditem

%beginitem
2. 
Написати функцію, що повертає \textbf{новий список} з записаною в нього послідовністю чисел
$$\scalebox{1.2}{$(-1)^{k}C_{3k}^{k}$}, \quad k=2,\ldots,n$$
Використовувати **, вкладені цикли та виклики функцій (крім методів класу list) \textbf{заборонено}.

%enditem
\end{large}
%end
\newpage
%begin
%\begin {Huge}
{{27.11.2025 Контрольна робота \No 2 (частина 2), варіант  \No \textbf { 11 }\par}}

%\begin{center}Частина 2
%\end{center}
%\vspace{10mm}
%\end {Huge}

{
\begin {large}
Якість коду та економність обчислень оцінюється по всім завданням.

%\vspace{5mm}

%beginitem
1. Написати функцію, що для інтервалу $[a, b)$ знаходить кількість чисел, квадрати яких діляться на суму квадратів їх цифр. (Розглядаємо десятковий запис.)

%enditem

%beginitem
2. 
Написати функцію, що повертає \textbf{новий список} з записаною в нього послідовністю чисел
$$\scalebox{1.2}{$C_{n+2k}^{k}$}, \quad k=0,1,\ldots,n-2$$
Використовувати **, вкладені цикли та виклики функцій (крім методів класу list) \textbf{заборонено}.

%enditem
\end{large}
%end
\newpage
%begin
%\begin {Huge}
{{27.11.2025 Контрольна робота \No 2 (частина 2), варіант  \No \textbf { 12 }\par}}

%\begin{center}Частина 2
%\end{center}
%\vspace{10mm}
%\end {Huge}

{
\begin {large}
Якість коду та економність обчислень оцінюється по всім завданням.

%\vspace{5mm}

%beginitem
1. Написати функцію, що для списку цілих знаходить найдовший відрізок, в якому всі сусідні числа різні. Якщо таких відрізків кілька, то повернути перший (як індекс початку та довжину). Використовувати додаткові структури даних (у тому числі рядки) заборонено.

%enditem

%beginitem
2. 
Написати функцію, що повертає \textbf{новий список} з записаною в нього послідовністю чисел
$$\scalebox{1.2}{$C_{n+2k}^{k}$}, \quad k=1,\ldots,n$$
Використовувати **, вкладені цикли та виклики функцій (крім методів класу list) \textbf{заборонено}.

%enditem
\end{large}
%end
\newpage
%begin
%\begin {Huge}
{{27.11.2025 Контрольна робота \No 2 (частина 2), варіант  \No \textbf { 13 }\par}}

%\begin{center}Частина 2
%\end{center}
%\vspace{10mm}
%\end {Huge}

{
\begin {large}
Якість коду та економність обчислень оцінюється по всім завданням.

%\vspace{5mm}

%beginitem
1. Написати функцію, що знаходить найменше число з списку цілих, у десятковому запису якого нема цифри 5. Повертає індекс останнього входження цього числа або $-1$ (за відсутності). Використовувати додаткові структури даних (у тому числі рядки) заборонено.

%enditem

%beginitem
2. 
Написати функцію, що повертає \textbf{новий список} з записаною в нього послідовністю чисел
$$\scalebox{1.2}{$C_{2k-1}^{k}$}, \quad k=2,\ldots,n$$
Використовувати **, вкладені цикли та виклики функцій (крім методів класу list) \textbf{заборонено}.

%enditem
\end{large}
%end
\newpage
%begin
%\begin {Huge}
{{27.11.2025 Контрольна робота \No 2 (частина 2), варіант  \No \textbf { 14 }\par}}

%\begin{center}Частина 2
%\end{center}
%\vspace{10mm}
%\end {Huge}

{
\begin {large}
Якість коду та економність обчислень оцінюється по всім завданням.

%\vspace{5mm}

%beginitem
1. Написати функцію, що знаходить найменше число з списку цілих, у десятковому запису якого нема цифри 7. Повертає індекс першого входження цього числа або $-1$ (за відсутності). Використовувати додаткові структури даних (у тому числі рядки) заборонено.

%enditem

%beginitem
2. 
Написати функцію, що повертає \textbf{новий список} з записаною в нього послідовністю чисел
$$\scalebox{1.2}{$(-1)^{k}C_{2k+1}^{k}$}, \quad k=2,\ldots,n$$
Використовувати **, вкладені цикли та виклики функцій (крім методів класу list) \textbf{заборонено}.

%enditem
\end{large}
%end
\newpage
%begin
%\begin {Huge}
{{27.11.2025 Контрольна робота \No 2 (частина 2), варіант  \No \textbf { 15 }\par}}

%\begin{center}Частина 2
%\end{center}
%\vspace{10mm}
%\end {Huge}

{
\begin {large}
Якість коду та економність обчислень оцінюється по всім завданням.

%\vspace{5mm}

%beginitem
1. Написати функцію, що для списку цілих знаходить найдовший відрізок, в якому всі сусідні числа відрізняються якнайменш на 3. Якщо таких відрізків кілька, то повернути перший (як індекс початку та довжину). Використовувати додаткові структури даних (у тому числі рядки) заборонено.

%enditem

%beginitem
2. 
Написати функцію, що повертає \textbf{новий список} з записаною в нього послідовністю чисел
$$\scalebox{1.2}{$(-1)^{k}C_{n+k}^{2k}$}, \quad k=1,\ldots,n$$
Використовувати **, вкладені цикли та виклики функцій (крім методів класу list) \textbf{заборонено}.

%enditem
\end{large}
%end
\newpage
%begin
%\begin {Huge}
{{27.11.2025 Контрольна робота \No 2 (частина 2), варіант  \No \textbf { 16 }\par}}

%\begin{center}Частина 2
%\end{center}
%\vspace{10mm}
%\end {Huge}

{
\begin {large}
Якість коду та економність обчислень оцінюється по всім завданням.

%\vspace{5mm}

%beginitem
1. Написати функцію, що знаходить найбільше число з списку цілих, у десятковому запису якого нема цифри 3. Повертає індекс останнього входження цього числа або $-1$ (за відсутності). Використовувати додаткові структури даних (у тому числі рядки) заборонено.

%enditem

%beginitem
2. 
Написати функцію, що повертає \textbf{новий список} з записаною в нього послідовністю чисел
$$\scalebox{1.2}{$C_{n+2k}^{k+1}$}, \quad k=1,\ldots,n$$
Використовувати **, вкладені цикли та виклики функцій (крім методів класу list) \textbf{заборонено}.

%enditem
\end{large}
%end
\newpage
%begin
%\begin {Huge}
{{27.11.2025 Контрольна робота \No 2 (частина 2), варіант  \No \textbf { 17 }\par}}

%\begin{center}Частина 2
%\end{center}
%\vspace{10mm}
%\end {Huge}

{
\begin {large}
Якість коду та економність обчислень оцінюється по всім завданням.

%\vspace{5mm}

%beginitem
1. Написати функцію, що знаходить найбільше число з списку цілих, у десятковому запису якого нема цифри 4. Повертає індекс першого входження цього числа або $-1$ (за відсутності). Використовувати додаткові структури даних (у тому числі рядки) заборонено.

%enditem

%beginitem
2. 
Написати функцію, що повертає \textbf{новий список} з записаною в нього послідовністю чисел
$$\scalebox{1.2}{$(-1)^{k}C_{n+k}^{2k}$}, \quad k=0,1,\ldots,n$$
Використовувати **, вкладені цикли та виклики функцій (крім методів класу list) \textbf{заборонено}.

%enditem
\end{large}
%end
\newpage
%begin
%\begin {Huge}
{{27.11.2025 Контрольна робота \No 2 (частина 2), варіант  \No \textbf { 18 }\par}}

%\begin{center}Частина 2
%\end{center}
%\vspace{10mm}
%\end {Huge}

{
\begin {large}
Якість коду та економність обчислень оцінюється по всім завданням.

%\vspace{5mm}

%beginitem
1. Написати функцію, що в інтервалі $[a, b)$ знаходить ціле число, що ділиться на найбільший степінь числа 3. Якщо таких чисел декілька, повернути найменше.

%enditem

%beginitem
2. 
Написати функцію, що повертає \textbf{новий список} з записаною в нього послідовністю чисел
$$\scalebox{1.2}{$(-1)^{k}C_{2n+k}^{n+k}$}, \quad k=1,\ldots,n$$
Використовувати **, вкладені цикли та виклики функцій (крім методів класу list) \textbf{заборонено}.

%enditem
\end{large}
%end
\newpage
%begin
%\begin {Huge}
{{27.11.2025 Контрольна робота \No 2 (частина 2), варіант  \No \textbf { 19 }\par}}

%\begin{center}Частина 2
%\end{center}
%\vspace{10mm}
%\end {Huge}

{
\begin {large}
Якість коду та економність обчислень оцінюється по всім завданням.

%\vspace{5mm}

%beginitem
1. Написати функцію, що для інтервалу $[a, b)$ знаходить найближчу пару чисел з цього інтервалу, що мають від трьох до п'яти (включно) простих дільників.

%enditem

%beginitem
2. 
Написати функцію, що повертає \textbf{новий список} з записаною в нього послідовністю чисел
$$\scalebox{1.2}{$(-1)^{k}C_{n+k}^{k+1}$}, \quad k=0,1,\ldots,n$$
Використовувати **, вкладені цикли та виклики функцій (крім методів класу list) \textbf{заборонено}.

%enditem
\end{large}
%end
\newpage
%begin
%\begin {Huge}
{{27.11.2025 Контрольна робота \No 2 (частина 2), варіант  \No \textbf { 20 }\par}}

%\begin{center}Частина 2
%\end{center}
%\vspace{10mm}
%\end {Huge}

{
\begin {large}
Якість коду та економність обчислень оцінюється по всім завданням.

%\vspace{5mm}

%beginitem
1. Написати функцію, що для цілого числа знаходить кількість його простих дільників, що входять у його розклад не більш ніж у 5му степеню. \\ Наприклад, для числа $3^8{\cdot}7^9{\cdot}11{\cdot}23^{2}$ таких дільників два – $11$ та $23$, тому функція має повернути $2$.

%enditem

%beginitem
2. 
Написати функцію, що повертає \textbf{новий список} з записаною в нього послідовністю чисел
$$\scalebox{1.2}{$(-1)^{k}C_{n+2k}^{k+1}$}, \quad k=1,\ldots,n$$
Використовувати **, вкладені цикли та виклики функцій (крім методів класу list) \textbf{заборонено}.

%enditem
\end{large}
%end
\newpage
%begin
%\begin {Huge}
{{27.11.2025 Контрольна робота \No 2 (частина 2), варіант  \No \textbf { 21 }\par}}

%\begin{center}Частина 2
%\end{center}
%\vspace{10mm}
%\end {Huge}

{
\begin {large}
Якість коду та економність обчислень оцінюється по всім завданням.

%\vspace{5mm}

%beginitem
1. Написати функцію, що для інтервалу $[a, b)$ знаходить найближчу пару чисел з цього інтервалу, що мають не менше трьох простих дільників та мають цифру 3 у своєму десятковому запису.

%enditem

%beginitem
2. 
Написати функцію, що повертає \textbf{новий список} з записаною в нього послідовністю чисел
$$\scalebox{1.2}{$(-1)^{k}C_{n+2k}^{k+1}$}, \quad k=0,1,\ldots,n$$
Використовувати **, вкладені цикли та виклики функцій (крім методів класу list) \textbf{заборонено}.

%enditem
\end{large}
%end
\newpage
%begin
%\begin {Huge}
{{27.11.2025 Контрольна робота \No 2 (частина 2), варіант  \No \textbf { 22 }\par}}

%\begin{center}Частина 2
%\end{center}
%\vspace{10mm}
%\end {Huge}

{
\begin {large}
Якість коду та економність обчислень оцінюється по всім завданням.

%\vspace{5mm}

%beginitem
1. Написати функцію, що повертає кількість входжень підрядка XYXYX у рядок. Наприклад, рядок ZXYXYXYXZSD містить два входження. Використовувати додаткові структури даних заборонено. Використовувати додаткові структури даних (у тому числі рядки) заборонено.

%enditem

%beginitem
2. 
Написати функцію, що повертає \textbf{новий список} з записаною в нього послідовністю чисел
$$\scalebox{1.2}{$(-1)^{k}C_{3k-1}^{k}$}, \quad k=2,\ldots,n$$
Використовувати **, вкладені цикли та виклики функцій (крім методів класу list) \textbf{заборонено}.

%enditem
\end{large}
%end
\newpage
%begin
%\begin {Huge}
{{27.11.2025 Контрольна робота \No 2 (частина 2), варіант  \No \textbf { 23 }\par}}

%\begin{center}Частина 2
%\end{center}
%\vspace{10mm}
%\end {Huge}

{
\begin {large}
Якість коду та економність обчислень оцінюється по всім завданням.

%\vspace{5mm}

%beginitem
1. Написати функцію, що для інтервалу $[a, b)$ знаходить найближчу пару чисел з цього інтервалу, що кожне з них має непарну кількість різних простих дільників. Якщо таких пар кілька, то цікавить пара з найбільшою сумою.

%enditem

%beginitem
2. 
Написати функцію, що повертає \textbf{новий список} з записаною в нього послідовністю чисел
$$\scalebox{1.2}{$(-1)^{k}C_{n+2k}^{k}$}, \quad k=1,\ldots,n$$
Використовувати **, вкладені цикли та виклики функцій (крім методів класу list) \textbf{заборонено}.

%enditem
\end{large}
%end
\newpage
%begin
%\begin {Huge}
{{27.11.2025 Контрольна робота \No 2 (частина 2), варіант  \No \textbf { 24 }\par}}

%\begin{center}Частина 2
%\end{center}
%\vspace{10mm}
%\end {Huge}

{
\begin {large}
Якість коду та економність обчислень оцінюється по всім завданням.

%\vspace{5mm}

%beginitem
1. Написати функцію, яка для цілого числа знаходить кількість його різних простих дільників, що мають у десятковому запису цифру 1.

%enditem

%beginitem
2. 
Написати функцію, що повертає \textbf{новий список} з записаною в нього послідовністю чисел
$$\scalebox{1.2}{$(-1)^{k}C_{2k-1}^{k}$}, \quad k=2,\ldots,n$$
Використовувати **, вкладені цикли та виклики функцій (крім методів класу list) \textbf{заборонено}.

%enditem
\end{large}
%end
\newpage
%begin
%\begin {Huge}
{{27.11.2025 Контрольна робота \No 2 (частина 2), варіант  \No \textbf { 25 }\par}}

%\begin{center}Частина 2
%\end{center}
%\vspace{10mm}
%\end {Huge}

{
\begin {large}
Якість коду та економність обчислень оцінюється по всім завданням.

%\vspace{5mm}

%beginitem
1. Написати функцію, що для списку цілих знаходить довжину найдовшого відрізку, в якому всі сусідні числа різні. Використовувати додаткові структури даних (у тому числі рядки) заборонено.

%enditem

%beginitem
2. 
Написати функцію, що повертає \textbf{новий список} з записаною в нього послідовністю чисел
$$\scalebox{1.2}{$C_{n+2k}^{k+1}$}, \quad k=0,1,\ldots,n$$
Використовувати **, вкладені цикли та виклики функцій (крім методів класу list) \textbf{заборонено}.

%enditem
\end{large}
%end
\newpage
%begin
%\begin {Huge}
{{27.11.2025 Контрольна робота \No 2 (частина 2), варіант  \No \textbf { 26 }\par}}

%\begin{center}Частина 2
%\end{center}
%\vspace{10mm}
%\end {Huge}

{
\begin {large}
Якість коду та економність обчислень оцінюється по всім завданням.

%\vspace{5mm}

%beginitem
1. Написати функцію, що для інтервалу $[a, b)$ знаходить найближчу пару чисел з цього інтервалу, щокожне з них має парну кількість різних простих дільників. Якщо таких пар кілька, то цікавить пара з найбільшою сумою.

%enditem

%beginitem
2. 
Написати функцію, що повертає \textbf{новий список} з записаною в нього послідовністю чисел
$$\scalebox{1.2}{$C_{n+k}^{2k}$}, \quad k=0,1,\ldots,n$$
Використовувати **, вкладені цикли та виклики функцій (крім методів класу list) \textbf{заборонено}.

%enditem
\end{large}
%end
\newpage
%begin
%\begin {Huge}
{{27.11.2025 Контрольна робота \No 2 (частина 2), варіант  \No \textbf { 27 }\par}}

%\begin{center}Частина 2
%\end{center}
%\vspace{10mm}
%\end {Huge}

{
\begin {large}
Якість коду та економність обчислень оцінюється по всім завданням.

%\vspace{5mm}

%beginitem
1. Написати функцію, що повертає число з списку цілих, що має парну кількість різних простих дільників. Якщо таких чисел кілька, то повернути найбільше. Використовувати додаткові структури даних (у тому числі рядки) заборонено.

%enditem

%beginitem
2. 
Написати функцію, що повертає \textbf{новий список} з записаною в нього послідовністю чисел
$$\scalebox{1.2}{$C_{3k}^{k}$}, \quad k=2,\ldots,n$$
Використовувати **, вкладені цикли та виклики функцій (крім методів класу list) \textbf{заборонено}.

%enditem
\end{large}
%end
\newpage
%begin
%\begin {Huge}
{{27.11.2025 Контрольна робота \No 2 (частина 2), варіант  \No \textbf { 28 }\par}}

%\begin{center}Частина 2
%\end{center}
%\vspace{10mm}
%\end {Huge}

{
\begin {large}
Якість коду та економність обчислень оцінюється по всім завданням.

%\vspace{5mm}

%beginitem
1. Написати функцію, що для цілого числа знаходить кількість його простих дільників, що входять у його розклад рівно у другому степеню. \\ Наприклад, для числа $3^2{\cdot}7^9{\cdot}11{\cdot}23^{2}$ таких дільників два – $3$ та $23$, тому функція має повернути $2$.

%enditem

%beginitem
2. 
Написати функцію, що повертає \textbf{новий список} з записаною в нього послідовністю чисел
$$\scalebox{1.2}{$(-1)^{k}C_{n+k}^{k}$}, \quad k=1,\ldots,n-2$$
Використовувати **, вкладені цикли та виклики функцій (крім методів класу list) \textbf{заборонено}.

%enditem
\end{large}
%end
\newpage
%begin
%\begin {Huge}
{{27.11.2025 Контрольна робота \No 2 (частина 2), варіант  \No \textbf { 29 }\par}}

%\begin{center}Частина 2
%\end{center}
%\vspace{10mm}
%\end {Huge}

{
\begin {large}
Якість коду та економність обчислень оцінюється по всім завданням.

%\vspace{5mm}

%beginitem
1. Написати функцію, що в інтервалі $[a, b)$ знаходить ціле число, що ділиться на найбільший степінь числа 3. Якщо таких чисел декілька, повернути найбільше.

%enditem

%beginitem
2. 
Написати функцію, що повертає \textbf{новий список} з записаною в нього послідовністю чисел
$$\scalebox{1.2}{$C_{2n+k+1}^{n+k}$}, \quad k=1,\ldots,n$$
Використовувати **, вкладені цикли та виклики функцій (крім методів класу list) \textbf{заборонено}.

%enditem
\end{large}
%end
\newpage
%begin
%\begin {Huge}
{{27.11.2025 Контрольна робота \No 2 (частина 2), варіант  \No \textbf { 30 }\par}}

%\begin{center}Частина 2
%\end{center}
%\vspace{10mm}
%\end {Huge}

{
\begin {large}
Якість коду та економність обчислень оцінюється по всім завданням.

%\vspace{5mm}

%beginitem
1. Написати функцію, що для цілого числа знаходить простий дільник, який входить у розклад числа в найбільшому степені. Якщо таких кілька, то повернути найменший. \\ Наприклад, для числа $2^{9}{\cdot}3^5{\cdot}7^9{\cdot}11$ функція має повернути $2$.

%enditem

%beginitem
2. 
Написати функцію, що повертає \textbf{новий список} з записаною в нього послідовністю чисел
$$\scalebox{1.2}{$(-1)^{k}C_{n+k}^{k}$}, \quad k=0,1,\ldots,n$$
Використовувати **, вкладені цикли та виклики функцій (крім методів класу list) \textbf{заборонено}.

%enditem
\end{large}
%end
\newpage
%begin
%\begin {Huge}
{{27.11.2025 Контрольна робота \No 2 (частина 2), варіант  \No \textbf { 31 }\par}}

%\begin{center}Частина 2
%\end{center}
%\vspace{10mm}
%\end {Huge}

{
\begin {large}
Якість коду та економність обчислень оцінюється по всім завданням.

%\vspace{5mm}

%beginitem
1. Написати функцію, що для списку цілих знаходить довжину найдовшого відрізку, що складається з чисел, у сімковому запису якого нема цифри 2. За відсутності має повертати $0$. Використовувати додаткові структури даних (у тому числі рядки) заборонено.

%enditem

%beginitem
2. 
Написати функцію, що повертає \textbf{новий список} з записаною в нього послідовністю чисел
$$\scalebox{1.2}{$(-1)^{k}C_{2n+k+1}^{n+k}$}, \quad k=0,1,\ldots,n$$
Використовувати **, вкладені цикли та виклики функцій (крім методів класу list) \textbf{заборонено}.

%enditem
\end{large}
%end
\newpage
%begin
%\begin {Huge}
{{27.11.2025 Контрольна робота \No 2 (частина 2), варіант  \No \textbf { 32 }\par}}

%\begin{center}Частина 2
%\end{center}
%\vspace{10mm}
%\end {Huge}

{
\begin {large}
Якість коду та економність обчислень оцінюється по всім завданням.

%\vspace{5mm}

%beginitem
1. Написати функцію, що для цілого числа знаходить суму його простих дільників, що входять у його розклад не більш ніж у 4му степеню. \\ Наприклад, для числа $3^5{\cdot}7^9{\cdot}11{\cdot}23^{2}$ таких дільників два – $11$ та $23$, тому функція має повернути $34$.

%enditem

%beginitem
2. 
Написати функцію, що повертає \textbf{новий список} з записаною в нього послідовністю чисел
$$\scalebox{1.2}{$(-1)^{k}C_{n+k}^{k+1}$}, \quad k=1,\ldots,n$$
Використовувати **, вкладені цикли та виклики функцій (крім методів класу list) \textbf{заборонено}.

%enditem
\end{large}
%end
\newpage
%begin
%\begin {Huge}
{{27.11.2025 Контрольна робота \No 2 (частина 2), варіант  \No \textbf { 33 }\par}}

%\begin{center}Частина 2
%\end{center}
%\vspace{10mm}
%\end {Huge}

{
\begin {large}
Якість коду та економність обчислень оцінюється по всім завданням.

%\vspace{5mm}

%beginitem
1. Написати функцію, що повертає найменше число з списку цілих, що має не більше трьох різних простих дільників. За відсутності має повертати None. Використовувати додаткові структури даних (у тому числі рядки) заборонено.

%enditem

%beginitem
2. 
Написати функцію, що повертає \textbf{новий список} з записаною в нього послідовністю чисел
$$\scalebox{1.2}{$C_{2n+k+1}^{n+k}$}, \quad k=0,1,\ldots,n$$
Використовувати **, вкладені цикли та виклики функцій (крім методів класу list) \textbf{заборонено}.

%enditem
\end{large}
%end
\newpage
%begin
%\begin {Huge}
{{27.11.2025 Контрольна робота \No 2 (частина 2), варіант  \No \textbf { 34 }\par}}

%\begin{center}Частина 2
%\end{center}
%\vspace{10mm}
%\end {Huge}

{
\begin {large}
Якість коду та економність обчислень оцінюється по всім завданням.

%\vspace{5mm}

%beginitem
1. Написати функцію, що повертає кількість входжень підрядка XYZYX у рядок. Наприклад, рядок ZXYZYXYZYXD містить два входження. Використовувати додаткові структури даних (у тому числі рядки) заборонено.

%enditem

%beginitem
2. 
Написати функцію, що повертає \textbf{новий список} з записаною в нього послідовністю чисел
$$\scalebox{1.2}{$C_{n+k}^{k+1}$}, \quad k=1,\ldots,n$$
Використовувати **, вкладені цикли та виклики функцій (крім методів класу list) \textbf{заборонено}.

%enditem
\end{large}
%end
\newpage
%begin
%\begin {Huge}
{{27.11.2025 Контрольна робота \No 2 (частина 2), варіант  \No \textbf { 35 }\par}}

%\begin{center}Частина 2
%\end{center}
%\vspace{10mm}
%\end {Huge}

{
\begin {large}
Якість коду та економність обчислень оцінюється по всім завданням.

%\vspace{5mm}

%beginitem
1. Написати функцію, що в інтервалі $[a, b)$ знаходить ціле число, що ділиться на найбільший степінь числа 3. Якщо таких чисел декілька, повернути найбільше.

%enditem

%beginitem
2. 
Написати функцію, що повертає \textbf{новий список} з записаною в нього послідовністю чисел
$$\scalebox{1.2}{$(-1)^{k}C_{n+k}^{2k-1}$}, \quad k=1,\ldots,n$$
Використовувати **, вкладені цикли та виклики функцій (крім методів класу list) \textbf{заборонено}.

%enditem
\end{large}
%end
\newpage
%begin
%\begin {Huge}
{{27.11.2025 Контрольна робота \No 2 (частина 2), варіант  \No \textbf { 36 }\par}}

%\begin{center}Частина 2
%\end{center}
%\vspace{10mm}
%\end {Huge}

{
\begin {large}
Якість коду та економність обчислень оцінюється по всім завданням.

%\vspace{5mm}

%beginitem
1. Написати функцію, що повертає число з списку цілих, що має найменшу кількість різних простих дільників. Якщо таких чисел кілька, то повернути найбільше. Використовувати додаткові структури даних (у тому числі рядки) заборонено.

%enditem

%beginitem
2. 
Написати функцію, що повертає \textbf{новий список} з записаною в нього послідовністю чисел
$$\scalebox{1.2}{$C_{n+k}^{k}$}, \quad k=0,1,\ldots,n-2$$
Використовувати **, вкладені цикли та виклики функцій (крім методів класу list) \textbf{заборонено}.

%enditem
\end{large}
%end
\newpage
%begin
%\begin {Huge}
{{27.11.2025 Контрольна робота \No 2 (частина 2), варіант  \No \textbf { 37 }\par}}

%\begin{center}Частина 2
%\end{center}
%\vspace{10mm}
%\end {Huge}

{
\begin {large}
Якість коду та економність обчислень оцінюється по всім завданням.

%\vspace{5mm}

%beginitem
1. Написати функцію, що для цілого числа знаходить простий дільник, який входить у розклад числа в найбільшому степені. Якщо таких кілька, то повернути найбільший. \\ Наприклад, для числа $2^{9}{\cdot}3^5{\cdot}7^9{\cdot}11$ функція має повернути $7$.

%enditem

%beginitem
2. 
Написати функцію, що повертає \textbf{новий список} з записаною в нього послідовністю чисел
$$\scalebox{1.2}{$(-1)^{k}\frac{C_{2k}^k}{k+1}$}, \quad k=0,1,\ldots,n-2$$
Використовувати **, вкладені цикли та виклики функцій (крім методів класу list) \textbf{заборонено}.

%enditem
\end{large}
%end
\newpage
%begin
%\begin {Huge}
{{27.11.2025 Контрольна робота \No 2 (частина 2), варіант  \No \textbf { 38 }\par}}

%\begin{center}Частина 2
%\end{center}
%\vspace{10mm}
%\end {Huge}

{
\begin {large}
Якість коду та економність обчислень оцінюється по всім завданням.

%\vspace{5mm}

%beginitem
1. Написати функцію, що повертає найбільше число з списку цілих, що має рівно три різних простих дільники. За відсутності має повертати None. Використовувати додаткові структури даних (у тому числі рядки) заборонено.

%enditem

%beginitem
2. 
Написати функцію, що повертає \textbf{новий список} з записаною в нього послідовністю чисел
$$\scalebox{1.2}{$C_{n+k}^{2k-1}$}, \quad k=1,\ldots,n$$
Використовувати **, вкладені цикли та виклики функцій (крім методів класу list) \textbf{заборонено}.

%enditem
\end{large}
%end
\newpage
%begin
%\begin {Huge}
{{27.11.2025 Контрольна робота \No 2 (частина 2), варіант  \No \textbf { 39 }\par}}

%\begin{center}Частина 2
%\end{center}
%\vspace{10mm}
%\end {Huge}

{
\begin {large}
Якість коду та економність обчислень оцінюється по всім завданням.

%\vspace{5mm}

%beginitem
1. Написати функцію, що для списку цілих знаходить довжину найдовшого відрізку, в якому всі сусідні числа відрізняються якнайменш на 2. Використовувати додаткові структури даних (у тому числі рядки) заборонено.

%enditem

%beginitem
2. 
Написати функцію, що повертає \textbf{новий список} з записаною в нього послідовністю чисел
$$\scalebox{1.2}{$\frac{C_{2k}^k}{k+1}$}, \quad k=0,1,\ldots,n$$
Використовувати **, вкладені цикли та виклики функцій (крім методів класу list) \textbf{заборонено}.

%enditem
\end{large}
%end
\newpage
%begin
%\begin {Huge}
{{27.11.2025 Контрольна робота \No 2 (частина 2), варіант  \No \textbf { 40 }\par}}

%\begin{center}Частина 2
%\end{center}
%\vspace{10mm}
%\end {Huge}

{
\begin {large}
Якість коду та економність обчислень оцінюється по всім завданням.

%\vspace{5mm}

%beginitem
1. Написати функцію, що для списку цілих знаходить найдовший відрізок, в якому всі сусідні числа різні. Якщо таких відрізків кілька, то повернути останній (як індекс початку та довжину). Використовувати додаткові структури даних (у тому числі рядки) заборонено.

%enditem

%beginitem
2. 
Написати функцію, що повертає \textbf{новий список} з записаною в нього послідовністю чисел
$$\scalebox{1.2}{$C_{n+k}^{k}$}, \quad k=1,\ldots,n$$
Використовувати **, вкладені цикли та виклики функцій (крім методів класу list) \textbf{заборонено}.

%enditem
\end{large}
%end
\newpage
\end{document}

